\begin{explanationbox}
\textbf{\textcolor{CExplanation}{一句话直觉}}

对于任意实数 $x$，正弦函数 $|\sin x|$ 的值总是被 $|x|$ 控制，当 $x$ 远离原点时 $|x|$ 更大，当 $x$ 接近原点时两者接近。

\medskip
\textbf{\textcolor{CExplanation}{核心拆解}}
\begin{description}[style=nextline, leftmargin=2.8em, labelsep=0.5em, itemsep=0.45em]
    \item[\textbf{要点一（绝对值比较）：}] 比较两个绝对值函数的大小关系。
    \item[\textbf{要点二（取等唯一）：}] 等号仅在 $x=0$ 时成立，其他位置严格小于。
\end{description}

\medskip
\textbf{\textcolor{CExplanation}{几何本质（弧长与弦长）}}

在单位圆中，$|x|$ 表示弧长（以弧度度量），$|\sin x|$ 表示正弦线长度。由于直线段最短，所以正弦线不超过弧长。

\medskip
\textbf{\textcolor{CExplanation}{代数意义（函数差单调）}}

考虑函数 $f(x)=x-\sin x$（当 $x\ge 0$），通过求导可得 $f'(x)=1-\cos x\ge 0$，故 $f(x)$ 单调不减，从而 $x\ge \sin x$；再结合奇偶性扩展到整个实数轴。

\medskip
\textbf{\textcolor{CExplanation}{考点价值（放缩工具）}}
\begin{itemize}[leftmargin=*, itemsep=0.5em, topsep=0.4em]
    \item \textbf{考法一（直接放缩）：} 在不等式证明中直接使用 $|\sin x|\le |x|$ 进行放缩。
    \item \textbf{考法二（极限估计）：} 用于计算极限 $\lim_{x\to 0}\frac{\sin x}{x}=1$ 的相关证明。
\end{itemize}

\medskip
\textbf{\textcolor{CExplanation}{顿悟点}}

\textbf{当 $x$ 很小时，$|\sin x|$ 与 $|x|$ 几乎相等，这解释了为什么在零点附近正弦函数可用线性近似。}

\medskip
\textbf{\textcolor{CExplanation}{使用场景}}
\begin{itemize}[leftmargin=*, itemsep=0.5em, topsep=0.4em]
    \item \textbf{场景一（证明不等式）：} 在涉及三角函数的不等式证明中，作为基础放缩依据。
    \item \textbf{场景二（参数范围）：} 用于确定含 $\sin x$ 的表达式中参数的取值范围。
\end{itemize}
\end{explanationbox}
